\section{Boolean encoding and preservation of candidates}\label{sec:encoding}
A formula in conjunctive normal form (CNF) is an AND of clauses, each
clause being an OR of literals. The encoding is built in three layers:
choose one answer at each observation, require a correct answer on every
proper coloring, then impose the selected branch counts and diagonal
answers. The first two layers describe the game exactly. The last uses
auxiliary variables, for which we need a preservation statement: a
hypothetical winning strategy in the branch must extend to a satisfying
assignment of the actual saved CNF.
\subsection{The base formula}
For a vertex $i\in\{1,\ldots,5\}$ and colors $a,b,g\in Q$ with
$g\notin\{a,b\}$,
introduce a Boolean variable
\[
x_{i,a,b,g}\quad\text{meaning}\quad f_i(a,b)=g.
\]
At each observation $(i,a,b)$ impose the clauses
\begin{align}
&\bigvee_{g\notin\{a,b\}}x_{i,a,b,g},\label{eq:at-least-one}\\
&\neg x_{i,a,b,g}\lor\neg x_{i,a,b,h}
\quad(g<h,\ g,h\notin\{a,b\}).\label{eq:at-most-one}
\end{align}
For every $c\in\calC_5$, add the coverage clause
\begin{equation}\label{eq:coverage-clause}
\bigvee_{i=1}^5x_{i,c_{i-1},c_{i+1},c_i}.
\end{equation}
All its variables exist because $c$ is proper. Let $\Phi$ be the
conjunction of these clauses.

For example, the proper coloring $(5,1,5,2,3)$ contributes the clause
\[
x_{1,3,1,5}\lor x_{2,5,5,1}\lor x_{3,1,2,5}
\lor x_{4,5,3,2}\lor x_{5,2,5,3}.
\]
Its five alternatives say exactly that player 1, 2, 3, 4 or 5 guesses
correctly on this coloring. The diagonal observation $(5,5)$ at player 2
is legitimate because only adjacent vertices must have different colors.

\begin{lemma}[Exactness of the base encoding]\label{lem:base}
The satisfying assignments of $\Phi$ are in bijection with normalized
five-color winning strategies.
\end{lemma}
\begin{proof}
Clauses~\eqref{eq:at-least-one} and~\eqref{eq:at-most-one} select exactly
one allowed answer at each observation, determining a unique normalized
strategy. Clause~\eqref{eq:coverage-clause} holds precisely when at least
one player guesses correctly on its indexed configuration. There is one
such clause for every member of $\calC_5$, so their conjunction is exactly
\eqref{eq:winning}. Conversely, a normalized winning strategy sets its
selected answer variables to true and satisfies all these clauses.
\end{proof}

For each player there are five equal-color observations with four allowed
answers and twenty unequal-color observations with three allowed answers.
The formula therefore has
\[
5(5\cdot4+20\cdot3)=400
\]
variables. The number of selection clauses is
\[
5\left[5\left(1+\binom42\right)+20\left(1+\binom32\right)\right]=575.
\]
Together with the 1020 coverage clauses, this gives 1595 clauses. In
particular, equal-color neighbor observations and colorings using fewer
than five colors have not been omitted.

\subsection{The count-branch formula}
Fix $a\in\calR_n$, and write $r_i=a_j$ when $i=1+2j$ modulo~5.
The intended additional constraints are
\begin{equation}\label{eq:cardinality}
\sum_{u,v\in\{1,2,3,4\}}x_{i,u,v,5}=r_i
\quad(1\leq i\leq5),\qquad x_{1,5,5,1}=1.
\end{equation}
They fix only five counts and one permissible normalization. In
particular, no positions of the $r_i$ selected entries and no other answer
values are fixed.

The recorded implementation uses PySAT sequential-counter components
for the five equalities~\cite{pysat}. Each component has exactly the
16 listed input variables and its own auxiliary variables. Denote the
actual generated component by $K_{i,r_i}(X_i,Y_i)$, where
$X_i=(x_{i,u,v,5})_{u,v\in\{1,2,3,4\}}$ and the $Y_i$ are disjoint.
The actual parent formula is
\begin{equation}\label{eq:parent}
F_a=\Phi\land x_{1,5,5,1}\land\bigwedge_{i=1}^5K_{i,r_i}(X_i,Y_i).
\end{equation}
For a refined parent, the formula at pattern $t\in\calT$, where $\calT$
is defined by~\eqref{eq:rgs}, is exactly
\begin{equation}\label{eq:child}
F_{a,t}=F_a\land\bigwedge_{i=1}^5x_{i,5,5,t_i}.
\end{equation}
The first of these five units duplicates the normalization in $F_a$ and
is harmless.

For the weight-17 example $a=(2,4,4,2,5)$, the counter targets are
$(r_1,r_2,r_3,r_4,r_5)=(2,2,4,5,4)$. In a possible refinement such as
$t=(1,2,1,3,2)$, the answers at $(5,5)$ are fixed to those five values,
while every other answer remains subject only to the base formula and
counts. Fixing a count is therefore different from fixing a table.

\subsection{Why auxiliary encodings do not discard a true strategy}
The required direction for an impossibility proof is
\begin{equation}\label{eq:extension}
\left(\sum X_i=r_i\right)
\ \Longrightarrow\ \exists Y_i\ K_{i,r_i}(X_i,Y_i).
\end{equation}
Accidentally removing an intended input could invalidate an exclusion
proof. Admitting extra inputs makes exclusion harder, but a refutation
would still exclude every intended input. Thus the reverse implication
is not needed. The audit checks~\eqref{eq:extension} on every intended
input of each actual component type, rather than relying only on a
library's claim of exactness. We do not assert a bijection between
strategies and auxiliary assignments; Lemma~\ref{lem:base} concerns only
the base formula.

For this check, the 16 input variables are renamed $1,\ldots,16$, and the
component's auxiliary variables are renamed consecutively from~17. A
component type is identified by its count and its complete normalized
clause list, not just by the count. After a particular input assignment
has been substituted, each remaining clause has at most one positive
auxiliary literal. This Horn property is checked from the actual clauses.

Under substitution, a clause with a true input literal is discarded;
false input literals are removed from the remaining clauses. A remaining
clause $\neg y_1\lor\cdots\lor\neg y_k\lor y$ is the rule
$y_1\land\cdots\land y_k\Rightarrow y$. With no positive literal its
head is $\bot$; with no negative literals it is a fact. This also covers
unit clauses and empty residual clauses.

\begin{lemma}[Auxiliary extension check]\label{lem:horn}
For a fixed input assignment, suppose the residual auxiliary formula is
Horn. Start with no auxiliary variable true. Repeatedly apply every rule
\[
y_{j_1}\land\cdots\land y_{j_k}\ \Longrightarrow\ y_\ell
\]
whose premises are true. Treat a clause with no positive literal as the
same rule with conclusion $\bot$. If closure is reached without deriving
$\bot$, setting exactly the derived variables true satisfies the residual
formula.
\end{lemma}
\begin{proof}
There are finitely many variables, and every successful application adds
a previously false variable, so the process terminates. At closure, a
rule with all premises true has its head true. A rule with head $\bot$
cannot have all premises true. Hence every residual Horn clause is
satisfied. Clauses already satisfied by the fixed input remain satisfied,
so this assignment also satisfies the original component.
\end{proof}

The implementation additionally evaluates every original component clause
on the constructed assignment. For count $r$, it enumerates every
$r$-element subset of the 16 input positions, with no sampling. It checks
the variable ranges and clause intervals to ensure that no component uses
another component's auxiliary variables or unrelated strategy variables.
Consequently, the five separate auxiliary assignments can be combined.

\begin{proposition}[Preservation by a branch input]\label{prop:input}
Suppose the base clauses, variable lists, extra units, and all properties
\eqref{eq:extension} have been checked for a branch. Every normalized
winning strategy in that branch extends to a satisfying assignment of
its actual formula~\eqref{eq:parent} or~\eqref{eq:child}.
\end{proposition}
\begin{proof}
Lemma~\ref{lem:base} supplies a satisfying assignment of $\Phi$.
The strategy satisfies the branch counts and the normalization, so
\eqref{eq:extension} supplies an auxiliary assignment for each component.
The auxiliary sets are disjoint, so their union satisfies all five
components simultaneously. In a refined branch the strategy also
satisfies its five diagonal units.
\end{proof}

This extension is chosen \emph{after} all vertex and color relabelings.
In particular, applying Lemma~\ref{lem:patterns} changes the base
assignment, and we then choose new auxiliary values for that assignment.
No assertion that a sequential-counter formula is syntactically
equivariant under color permutations is needed.

The input audit reconstructs~\eqref{eq:at-least-one}--\eqref{eq:coverage-clause}
directly from the game, compares the base clause multiset, checks all five
component boundaries, and confirms that the only other parent clause is
$x_{1,5,5,1}$. Every child is compared clause by clause with its parent
plus the five specified units. These checks address the mathematical input
to the solver, separately from any subsequent refutation.

